\documentclass[00-main.tex]{subfiles}

\begin{document}

\section{Practical Considerations}\label{sec:practical-consideration}
To contextualize the computational cost and parameter recovery of the proposed pseudo-MLE, we compare it to a Bayesian Markov Chain Monte Carlo (MCMC) estimator from the \emph{HSSM} software \citep{fengler_hssm_2026}, the successor to the widely used \emph{HDDM} software \citep{wiecki_hddm_2013}. Both methods were applied to the same simulated, correctly specified sequence of choices and reaction times with parameters $\bm\theta = \{a = 0.63, t_0 = 0.435, v = 2.23, z = 0.008\}$ \citep{matzke_psychological_2009}, $n=\{500, 1000, 5000, 10000\}$ trials, and $b=100$ replications. Likewise, both methods were initialized at the same starting parameter values and constrained to a single CPU core per replication while parallelizing across replications on an 11-core macOS machine; otherwise, default settings were used. \emph{HSSM} (v0.3.0) defaults to 1000 tuning iterations, 1000 posterior draws, and 2 independent chains run sequentially. We compare runtime in seconds (averaged across replications), and root mean squared error (RMSE) of parameter and standard error estimates:
\begin{equation}
    \text{RMSE}(\hat{\bm\theta}) = \sqrt{\frac{1}{b} \sum_{j=1}^b (\hat{\bm\theta}_j - \bm\theta)^2}, \quad \text{RMSE}(\widehat{\text{se}}(\hat{\bm\theta})) = \sqrt{\frac{1}{b} \sum_{j=1}^b \bigl(\widehat{\text{se}}(\hat{\bm\theta})_j - \text{se}(\hat{\bm\theta})\bigr)^2}.
\end{equation}
Here, $\hat{\bm\theta}_j$ is the point estimate for the $j$-th replication, $\widehat{\text{se}}(\hat{\bm\theta})_j$ is its estimated standard error (standard error for MLE, marginal posterior standard deviation for MCMC), and $\text{se}(\hat{\bm\theta})$ is the empirical standard deviation of the $b$ parameter estimates, serving as the true standard error.

Theoretically, parameters and standard errors are different when comparing MLE (Frequentist) and MCMC (Bayesian) methods, as the latter defines a prior on the parameter space. Additionally, MCMC estimates the full posterior distribution from which point estimates (marginal posterior means) and their standard errors (marginal posterior standard deviations) are calculated. Yet the influence of the prior diminishes at large samples, and asymptotically, the posterior converges to a normal distribution centered on the MLE with covariance $\mathscr{I}_{\bm\theta}^{-1}$ \citep[Theorem~10.1]{van_der_vaart_asymptotic_1998}; the same as the MLE under $iid$ trials and correct specification in Section~\ref{sec:properties}. This theorem links MLE and MCMC inference at large samples, and therefore we can compare the RMSE of parameter and standard error estimates of these methods.

Figure~\ref{fig:mle-mcmc} compares MLE and MCMC under $iid$ trials and correct specification, across four sequence lengths. MCMC takes $9$--$43\times$ longer, with the gap widening at larger $n$ as the per-sample likelihood evaluation cost compounds over many posterior samples. MLE is more computationally efficient because it uses first- and second-order likelihood information to converge directly to a point estimate, whereas MCMC must characterize the full posterior distribution through stochastic sampling. Despite this difference in computational cost, both methods recover parameters and their standard errors accurately, where both improve at larger $n$, and MLE performs as well or better than MCMC across most settings (except for $\text{RMSE}(\widehat{\text{se}}(\hat{v}))$ at $n=1000$ trials).

\begin{figure}[H]
    \centering
    \includegraphics[width=1\textwidth]{docs/manuscript/figures/fig07-mcmc.pdf}
    \caption{\emph{Comparison of MLE and MCMC.}}
    \subcaption*{
    Solid lines denote MLE, dashed lines denote MCMC, and different marker shapes distinguish the four parameters $a, t_0, v$, and $z$ across all panels.
    (A) Average runtime in seconds across replications as a function of sample size.
    (B) RMSE of the parameter estimates, in the original units of $a, t_0, v$, and $z$.
    (C) RMSE of the standard error estimates.
    }
    \label{fig:mle-mcmc}
\end{figure}

In context, both methods allow for covariate modeling and uncertainty quantification; however, MCMC as implemented assumes $iid$ trials and correct specification \citep{wiecki_hddm_2013, moens_recurrent_2018}. Therefore, standard errors are likely underestimated under temporal dependence and unmodeled parameter variability, while we show that pseudo-MLE robustly estimates standard errors in this setting (Figure~\ref{fig:sim_realistic}). That said, their MCMC approach has a lot to offer that we cannot. Most notably, the \emph{HSSM} software can approximate likelihoods when no analytical form exists, differentiable or not \citep{fengler_hssm_2026, fengler_beyond_2022, fengler_likelihood_2021}. This allows for a wide class of DDMs beyond the four-parameter model, including the seven-parameter model \citep{ratcliff_modeling_1998, ratcliff_estimating_2002}, collapsing-bound model \citep{drugowitsch_cost_2012}, L\'{e}vy-flight and Cauchy-flight models \citep{voss_sequential_2019}. Also, it can implement hierarchical structures to pool data across subjects or sessions, improving estimation accuracy when individual data is sparse. Therefore, while Bayesian MCMC offers extensive flexibility for complex and hierarchical generative models, our pseudo-MLE approach provides a highly scalable, explicitly robust alternative for fast inference and uncertainty quantification under temporal dependence and parameter variability.

\section{Conclusions}\label{sec:conclusions}
Standard DDM fitting methods treat trial sequences as independent and identically distributed ($iid$), an assumption routinely violated in naturalistic behavioral experiments where decisions are richly structured in time \citep{urai_structure_2025, moens_recurrent_2018}. This paper introduces a pseudo-MLE method that provides analytically tractable uncertainty quantification for DDM parameters while remaining valid under temporal dependence and unmodeled parameter variability, and extending to settings where parameters are explicitly modeled as functions of covariates.

\subsubsection*{Theoretical Contributions}
We establish that the pseudo-MLE converges to its pseudo-true coefficients and is asymptotically normal under strict stationarity and strong mixing. The key insight is that stationarity applies to the joint process of covariates and outcomes, not to the decision outcomes alone. By explicitly modeling structured nonstationarity --- such as session-level trends, circadian rhythms, and learning effects --- valid point estimates and uncertainties are recovered even in nonstationary examples, substantially broadening the scope of settings to which the method applies. The resulting sandwich-form covariance estimators \citep{newey_simple_1987, hansen_probability_2022} correct for the failure of the information matrix equality under possible misspecification, providing valid standard errors and confidence intervals even when the working model is an approximation of the true data-generating process.

\subsubsection*{Simulation Results}
Monte Carlo simulations confirm that under correct specification, with either constant or covariate-driven nonstationary parameters, point estimates are unbiased and covariances are accurately recovered, consistent with prior DDM parameter recovery studies \citep{boehm_estimating_2018, matzke_psychological_2009}. Under misspecification induced by unmodeled trial-to-trial parameter variability, non-robust covariance estimators substantially underestimate standard errors, while robust estimators largely mitigate these biases. Robust covariance estimation is therefore advisable whenever misspecification cannot be ruled out, which \citet{urai_structure_2025} cautions is the norm in decision-making experiments.

More broadly, explicitly modeling parameter variability via covariates is preferable to absorbing it into unstructured noise parameters, which carry no temporally structured information and conflate distinct underlying sources such as arousal, choice history, and learning \citep{duffy_disentangling_2025, urai_structure_2025, roy_efficient_2018, roy_extracting_2021}. A notable example is the seven-parameter DDM, where trial-to-trial variability in boundary height $a$ and starting point $z$ is difficult to disentangle, potentially leading to misattributed variability \citep{boehm_estimating_2018, lerche_model_2016, ratcliff_internal_2018}. Covariate modeling offers a more tractable alternative: by linking $a$ and $z$ separately to observed covariates, their respective sources of variability can be more cleanly separated. More generally, when structured variability can be tied to observable quantities, covariate modeling converts a source of inferential noise into interpretable signal about how decision states change over time.

\subsubsection*{Data Analysis Results}
Applied to empirical data from a single rat performing a visual motion discrimination task \citep{reinagel_speed_2013}, the method revealed systematic multi-timescale structure in DDM parameters invisible to standard analyses. At the between-day timescale, increasing decision boundaries and drift rate coefficients across the first 35 days could be a practice effect as often reported \citep{ashwood_mice_2022, mohammadi_identifying_2025, cochrane_multiple_2023}. Notably, parameters continued to vary well beyond this initial period, suggesting that decision states change even in well-trained animals. Comparing Figure~\ref{fig:rat-ddm}A,C further shows that explicitly modeling within-day nonstationary trends did not materially change the between-day point estimates, but it did sharpen inference by reducing uncertainty; across parameters, standard errors were smaller when trial-in-day covariates were included. Within a day, decision boundaries and non-decision times increased monotonically as the rat became satiated and fatigued, while markedly lower boundaries and faster responses in the first 20 trials might reflect heightened thirst at the start of the day. These findings show that even under mostly constant experimental conditions, DDM parameters varied substantially relative to their uncertainties, changes that cannot be inferred from point estimates alone and that methods ignoring trial order would attribute to noise. Persistent autocorrelations across trials and days further confirmed that sequential dependencies are present and should be accounted for in inference, consistent with history effects documented across rodent and primate decision making \citep{urai_choice_2019, busse_detection_2011, danskin_exponential_2023, frund_quantifying_2014, gao_sequential_2009}.

\subsubsection*{Limitations and Future Directions}
While misspecification-robust covariances reduce inferential errors, they cannot replace the explicit modeling of known sources of variability. For instance, history-dependent biases that drive sequential dependencies were not modeled here, as they were in \citep{nguyen_optimizing_2019, urai_choice_2019, braun_adaptive_2018, gupta_trial-history_2024}. Methodologically, the large-sample theory underlying our pseudo-MLE approach requires that the joint process of covariates and outcomes be strictly stationary and strongly mixing. These conditions are difficult to verify in practice, although our application to rat decision making suggests that realistic within-day nonstationarity did not substantively distort the between-day point estimates in this example (Figure~\ref{fig:rat-ddm}A,C); explicitly modeling it improved estimation precision. Empirically, applying the method to a single rat under fixed stimulus conditions was meant as a demonstration, and the results presented are specific to this example and not expected to generalize across different rats or task designs. We analyzed example covariates by way of illustration of the method, but did not attempt to exhaust the many potentially explanatory covariates in this dataset. Finally, because the current method estimates fixed coefficients within each fitted window, it characterizes continuous parameter shifts but cannot capture the abrupt, discrete strategy switches that latent-state models are specifically designed to detect \citep{ashwood_mice_2022, mohammadi_identifying_2025, kucharsky_hidden_2021}.

Future work could address these limitations in three primary ways. First, history-dependent biases and species-specific traits \citep{nguyen_different_2022, shevinsky_interaction_2019, li_dynamic_2024, pagan_individual_2025} are better addressed by incorporating them as covariates into the underlying model. Extending the method to large, multi-subject datasets would then enable the comparison of changing decision dynamics across different subjects, timescales, and species \citep{asadpour_rodents_2024, danskin_exponential_2023, nguyen_different_2022}. Second, testing additional covariates derived from external behavioral metrics, such as pose tracking \citep{mathis_deeplabcut_2018} or arousal signals \citep{murphy_pupil-linked_2014}, would give richer insights into how physical context shapes evidence accumulation. Linking these physical signals directly to DDM parameters would further disentangle structured variability explained by observed covariates and residual variability that requires robust inference. Finally, combining our pseudo-MLE approach with latent-state models would more fully capture the broad spectrum of temporal variability, such as hidden Markov models for discrete state switches \citep{ashwood_mice_2022, mohammadi_identifying_2025, kucharsky_hidden_2021}, autoregressive processes for gradual drift \citep{urai_choice_2019, vloeberghs_bayesian_2025}, transition models that combine discrete switches with continuous drift \citep{schumacher_neural_2023, schumacher_validation_2025}, or reinforcement learning for history effects \citep{pedersen_drift_2017, fengler_beyond_2022}. Together, these extensions could jointly account for structured nonstationarity, gradual parameter drift, and discrete state changes, more fully capturing the broad spectrum of temporal variability inherent in sequential decision making.

In conclusion, this work establishes a methodological foundation for studying dynamic decision making in naturalistic settings, in theory and applications. By relaxing the standard assumptions of correct specification and trial independence, and by grounding stationarity in the joint process of covariates and outcomes rather than in decision outcomes alone, the pseudo-MLE method enables principled and computationally efficient inference about sequential decision making in naturalistic experiments.
\end{document}
